A disseminação de arquiteturas distribuídas e ecossistemas nativos em nuvem baseados em
contêineres e microsserviços redefiniu os desafios de operação de software. Nessas arquiteturas, falhas podem decorrer tanto de componentes isolados quanto de
interações entre serviços e dependências assíncronas, o que amplia a complexidade do
diagnóstico. Nesse contexto, Kosińska et al.~\citeano{kosinska2023} discutem a importância
dos dados de telemetria gerados durante a execução das aplicações para a análise de
ambientes \emph{cloud-native}.

Fundamentada na Teoria de Controle clássica, a observabilidade é caracterizada pela
capacidade de um sistema permitir a identificação de seu estado interno a partir da
análise de suas saídas em um intervalo de tempo finito. Na engenharia de software, Alongi et al.~\citeano{alongi2022} discutem a observabilidade
como uma propriedade arquitetural relacionada à compreensão do comportamento do sistema.
Enquanto o monitoramento acompanha indicadores e condições previamente definidos,
a observabilidade favorece a investigação de comportamentos não antecipados a partir
dos dados disponíveis. Giamattei et al.~\citeano{giamattei2024} destacam que as
ferramentas operacionais possibilitam a coleta dessas informações, enquanto a observabilidade busca
organizá-las e relacioná-las de maneira que auxiliem na compreensão dos problemas e na
tomada de decisões pelas equipes de engenharia.

Essa visibilidade pode ser construída pela análise conjunta de três tipos de sinais
de telemetria:

\begin{samepage}
\subsubsection{Métricas}\label{sec:metricas}

Representam medições quantitativas coletadas ao longo do tempo, como
consumo de CPU, latência de resposta e taxas de erro. Essas informações permitem
acompanhar o desempenho e as condições gerais da infraestrutura e dos serviços.
\par
\end{samepage}

\begin{samepage}
\subsubsection{Logs}\label{sec:logs}

São registros cronológicos de eventos que armazenam informações sobre o
comportamento das aplicações, incluindo erros, exceções e mudanças de estado. Esses
registros ajudam a compreender o que aconteceu durante a execução de um sistema.
\par
\end{samepage}

\begin{samepage}
\subsubsection{Rastreamentos distribuídos (distributed tracing)}\label{sec:rastreamentos}

Permitem acompanhar o
percurso de uma requisição entre diferentes serviços e componentes de uma aplicação. Esse
acompanhamento facilita a identificação de falhas, atrasos e possíveis gargalos em
arquiteturas distribuídas.
\par
\end{samepage}

Giamattei et al.~\citeano{giamattei2024} e Kosińska et al.~\citeano{kosinska2023}
discutem ferramentas de coleta e análise de telemetria, como Prometheus, Jaeger, Fluentd
e Grafana. O OpenTelemetry oferece mecanismos abertos de instrumentação e integração
de telemetria que podem favorecer a interoperabilidade entre soluções e reduzir a
dependência de um único fornecedor (\emph{vendor lock-in}). Alongi et
al.~\citeano{alongi2022} discutem como abordagens como \emph{Event Sourcing} podem
ampliar a capacidade de auditoria e facilitar a investigação de incidentes. Entretanto, essa implementação exige considerar aspectos
como o esforço necessário para instrumentar as aplicações, o volume de dados armazenados e
os custos envolvidos na operação dessas soluções.

Quando essas necessidades técnicas de arquitetura e telemetria são levadas para o contexto
das instituições públicas, os desafios relacionados à confiabilidade dos sistemas passam a
envolver não apenas questões tecnológicas, mas também estruturas burocráticas, sistemas
legados e exigências legais.
